\section{Ordenação}
\begin{frame}[label=main]
	\begin{block}{Ordenar}\pause
		Processo de organizar um conjunto de objetos em uma ordem.
	\end{block}
	\vfill\pause
	\begin{itemize}
		\item lista telefônica, com os nomes das pessoas em ordem alfabética
		\item bancos, com os atendimentos por ordem de chegada
		\item hospitais, com os atendimentos em ordem de prioridade
		\item times, com os jogadores em ordem de técnica
		\item etc.
	\end{itemize}
	\vfill\pause
	A ordenação visa \alert{facilitar a recuperação de itens} do conjunto.
\end{frame}

\begin{frame}[label=main]
	\begin{block}{Notação}
		Algoritmos de ordenação trabalham sobre a \alert<1>{\emph{chave}} de registros ou estruturas [de um arquivo].
	\end{block}
	\pause
	\begin{itemize}
		\item nome
		\item matrícula
		\item id
		\item código
		\item etc.
	\end{itemize}
	\vfill
	Claro, podem existir outros componentes no registro.
\end{frame}

\lstset{morekeywords={chave_t, data_t}}
\begin{frame}[fragile,label=main]
	\begin{block}{Chave de Ordenação}
		Qualquer tipo de chave sobre o qual exista uma regra de ordenação bem definida.
	\end{block}
	\pause\vfill
	\begin{columns}[T]
		\column{.27\textwidth}
			\begin{lstlisting}[language=C]
					typedef int chave_t;
			\end{lstlisting}
		\column{.3\textwidth}
			\begin{lstlisting}[language=C]
					typedef char chave_t;
			\end{lstlisting}
			\vfill
			\begin{lstlisting}[title=Registro,language=C]
				typedef struct data_t{
				  chave_t key;
				  // etc.
				} data_t;
			\end{lstlisting}
		\column{.33\textwidth}
			\begin{lstlisting}[language=C]
				typedef struct chave_t;
			\end{lstlisting}
	\end{columns}
\end{frame}

\begin{frame}[label=main]
	\begin{block}{Métodos Estáveis}
		\uncover<2->{A ordem relativa dos itens com chaves iguais \emph{\alert<2>{não}} se altera durante a ordenação.}
	\end{block}
	\vfill
	\begin{columns}
		\column{.3\textwidth}\centering
			\emph{\small Vetor}\\
			\begin{tabular}{|l|}
				\hline
				\alert{C}arlos \\\hline
				\alert{A}lice \\\hline
				\alert{B}runo \\\hline
				\alert{A}fonso \\\hline
				\alert{B}ernardo \\\hline	
			\end{tabular}
			\pause
		\column{.3\textwidth}\centering
			\emph{\small Não Estável}\\
			\begin{tabular}{|l|}
				\hline
				\alert{A}fonso \\\hline
				\alert{A}lice \\\hline
				\alert{B}runo \\\hline	
				\alert{B}ernardo \\\hline
				\alert{C}arlos \\\hline
			\end{tabular}
			\pause
		\column{.3\textwidth}\centering
			\emph{\small Estável}\\
			\begin{tabular}{|l|}
				\hline
				\alert{A}lice \\\hline
				\alert{A}fonso \\\hline
				\alert{B}runo \\\hline
				\alert{B}ernardo \\\hline
				\alert{C}arlos \\\hline		
			\end{tabular}
	\end{columns}
	\vfill\pause
	Alguns dos métodos de ordenação mais eficientes não são estáveis, mas a estabilidade \emph{pode ser forçada}...
		%\subitem{Sedgewick (1988) sugeriu aumentar a chave antes da ordenação}
\end{frame}

\begin{frame}[label=main]
	\begin{center}
	\begin{minipage}{.25\textwidth}
		\begin{alertblock}{Ineficientes}
			\subitem{bogosort}
		\end{alertblock}
	\end{minipage}
	\vfill\pause
	\begin{columns}[T]
		\column{.3\textwidth}
			\begin{block}{Elementares}
				\begin{itemize}
					\item bubble sort
					\item inserção
					\item seleção
				\end{itemize}
			\end{block}
			% \pause
			% \begin{itemize}
			%	\item Adequados para pequenos arquivos.
			%	% \item Complexidade $O(n^2)$.
			%	\item Produzem programas pequenos.
			% \end{itemize}
			\pause
		\column{0\textwidth}
		\column{.3\textwidth}
			\begin{exampleblock}{Eficientes}
				\begin{itemize}
					\item quicksort
					\item shellsort
					\item heapsort
					\item mergesort
				\end{itemize}
			\end{exampleblock}
			% \pause
			% \begin{itemize}
			%	\item Adequados para arquivos maiores.
			%	% \item Complexidade $O(n\cdot\log(n))$.
			%	\item Produzem programas maiores.
			% \end{itemize}
	\end{columns}
	\end{center}
	\vfill\pause
	E outros: counting sort, radix sort, bucket sort, pancake sort, spaghetti sort, ...
\end{frame}